\chapter{可能形 — Forma Potencial}
\unidadheader{21}{可能形}{Forma Potencial}{Expresar lo que puedes o no puedes hacer}

\begin{tcolorbox}[vocabbox, title={📌 Vocabulario Nuevo}]
\begin{tabularx}{\linewidth}{llX}
\toprule
\textbf{Japonés} & \textbf{Español} & \textbf{Tipo} \\
\midrule
\vocabitem{泳ぐ}{およぐ}{nadar}{v. gr.1}
\vocabitem{歌う}{うたう}{cantar}{v. gr.1}
\vocabitem{弾く}{ひく}{tocar (instrumento)}{v. gr.1}
\vocabitem{運転する}{うんてんする}{conducir}{v. irr.}
\vocabitem{料理する}{りょうりする}{cocinar}{v. irr.}
\vocabitem{説明する}{せつめいする}{explicar}{v. irr.}
\vocabitem{翻訳する}{ほんやくする}{traducir}{v. irr.}
\vocabitem{〜が得意}{とくい}{ser bueno en ~}{expr.}
\vocabitem{なんとか}{—}{de alguna manera / más o menos}{adv.}
\vocabitem{全然}{ぜんぜん}{para nada (+ neg.)}{adv.}
\bottomrule
\end{tabularx}
\end{tcolorbox}

\begin{tcolorbox}[phrasebox, title={⚙️ Conjugación a forma potencial}]
\begin{tabular}{llll}
\toprule
\textbf{Grupo} & \textbf{Regla} & \textbf{Original} & \textbf{Potencial} \\
\midrule
1 (う→え段) & う→える & \ruby{書}{か}く & \ruby{書}{か}ける \\
2 (〜る → 〜られる) & る→られる & \ruby{食}{た}べる & \ruby{食}{た}べられる \\
3 (irregulares) & memorizar & する / くる & できる / こられる \\
\bottomrule
\end{tabular}
\end{tcolorbox}

\subsection*{🗣️ Frases Clave}

\frase{\ruby{日本語}{にほんご}が\ruby{話}{はな}せますか？}{¿Puedes hablar japonés?}
\frase{なんとか\ruby{話}{はな}せます。}{Más o menos puedo hablar.}
\frase{\ruby{全然}{ぜんぜん}\ruby{泳}{およ}げません。}{No puedo nadar para nada.}
\frase{ピアノが\ruby{弾}{ひ}けますか？}{¿Puedes tocar el piano?}

\subsection*{⚙️ Gramática}

\patron{Forma potencial: poder hacer ~}
  {V (potencial) + ます / V-ことができます}
  {\ej{\ruby{漢字}{かんじ}が\ruby{読}{よ}めます。}{Puedo leer kanji.}
  \ej{\ruby{車}{くるま}を\ruby{運転}{うんてん}することができません。}
    {No puedo conducir un coche.}
  \ej{スペイン\ruby{語}{ご}が\ruby{話}{はな}せますか？}
    {¿Puedes hablar español?}}

\begin{tcolorbox}[ejerciciobox, title={📝 Ejercicios Rápidos}]
\begin{enumerate}
  \item Convierte a potencial: \ruby{飲}{の}む / \ruby{見}{み}る / する
  \item Di 2 cosas que puedes hacer bien en japonés.
  \item Traduce: \ruby{子}{こ}どものとき、\ruby{木}{き}に\ruby{登}{のぼ}れましたか？
\end{enumerate}
\vspace{0.4em}
\textbf{Respuestas:}
1) \ruby{飲}{の}める / \ruby{見}{み}られる / できる\quad
2) (personal)\quad
3) Cuando eras niño/a, ¿podías subir a los árboles?
\end{tcolorbox}

\begin{tcolorbox}[kanjibox, title={🀄 Kanjis de la Unidad}]
\begin{tabularx}{\linewidth}{cllXX}
\toprule
\textbf{Kanji} & \textbf{On} & \textbf{Kun} & \textbf{Significado} & \textbf{Vocab} \\
\midrule
\kanjirow{泳}{えい}{およ}{nadar}{\ruby{泳}{およ}ぐ / \ruby{水泳}{すいえい}}
\kanjirow{歌}{か}{うた}{cantar / canción}{\ruby{歌}{うた}う / \ruby{歌手}{かしゅ}}
\kanjirow{運}{うん}{はこ}{transporte / suerte}{\ruby{運転}{うんてん} / \ruby{運動}{うんどう}}
\kanjirow{転}{てん}{ころ}{rodar / girar}{\ruby{運転}{うんてん} / \ruby{転ぶ}{ころぶ}}
\kanjirow{能}{のう}{—}{capacidad / habilidad}{\ruby{可能}{かのう} / \ruby{才能}{さいのう}}
\bottomrule
\end{tabularx}
\end{tcolorbox}
